($j$ is the common value of $j_1$ and $j_2$). The factor preceding the sum follows from normalization. All the coefficients in the sum must have the same absolute value, since every possible value of the particles' angular-momentum projection $m$ is equally probable. The pattern of alternating signs in (106.2) is readily obtained by writing the wave functions in spinor form. In spinor notation, the sum in (106.2) is the scalar (the system has zero total angular momentum!)
\begin{equation}
 \psi^{(1)}_{\lambda\mu\ldots}\psi^{(2)\lambda\mu\ldots},
 \tag{106.3}
\end{equation}
formed from two spinors of rank $2j$. Once this is recognized, the signs in (106.2) follow immediately from formula (57.3).
